\documentclass{article}
\usepackage{graphicx}
\usepackage[margin=1.5cm]{geometry}
\usepackage{amsmath}

\begin{document}
\twocolumn

\title{Monday Warm-Up: Unit 5, Electromagnetism}
\author{Prof. Jordan C. Hanson}

\maketitle

\section{Memory Bank}

\begin{itemize}
\item $U_{\rm C} = \frac{1}{2}Q^2C^{-1}$ ... Energy stored in a capacitor.
\item $C = \epsilon_0 A/d$ ... Parallel plate capacitance.
\item $U_{\rm C} = \frac{1}{2}LI^2$ ... Energy stored in an inductor.
\item $L = \mu_0 N^2 A/l$ ... Solenoidal inductance.
\item $c = f\lambda$ ... The relationship between speed of light, $c$, the frequency, $f$, and the wavelength, $\lambda$.
\item $c = 299,792,458$ m s$^{-1}$, or $c \approx 0.3$ m ns$^{-1}$.
\item $v = c/n$ ... The actual speed of light, $v$ in a material with index of refraction $n$.
\item $\bar{S} = E^2/(2c\mu_0)$ ... Average \textbf{intensity}, or W m$^{-2}$, of electromagnetic energy.
\item $\cos(x)\cos(y) = (1/2)(\cos(x+y) + \cos(x-y))$ ... Trigonometric identity for \textit{mixing} sinusoids.
\end{itemize}

\section{Electromagnetic Waves and \\ Energy}

\begin{enumerate}
\item (a) What is the intensity of a wave created by an E-field with magnitude 1000 V m$^{-1}$ at the source? (b) How do you think the previous result scales with distance from the source? \\ \vspace{2.5cm}
\item (a) What is the wavelength of a 6 GHz microwave in vacuum? (b) What is the frequency of a wave with wavelength 3 cm in vacuum? (d) If a wave has a wavelength of 10 cm in air, what is the wavelength in ice if the index of refraction is $n = 1.78$? \\ \vspace{3cm}
\end{enumerate}

\section{Signal Mixing and AM Radio}

\begin{enumerate}
\item (a) Suppose $f(t) = 4\cos(2\pi f_1 t)$ and $g(t) = 2\cos(2\pi f_2 t)$.  Using a formula from the Memory Bank, calculate the frequencies present in the product $f(t) g(t)$. (b) Suppose $f_1 = 1440$ kHz, and $f_2 = 560$ Hz.  What frequencies exist in $f(t) g(t)$? (c) Suppose an RLC resonator is built to take an audio signal at $f_2$ and mix it with a \textit{carrier frequency} at $f_1$.  The product is the \textit{radiated} across space.  Graph this radiated signal $f(t) g(t)$. \\ \vspace{3cm}
\end{enumerate}

\section{Geometric Optics}

\begin{enumerate}
\item Suppose a laser beam reflects from a mirror. (a) If the incident angle with respect to the normal direction of the surface is 30 degrees, what is the reflected angle? (b) Suppose two mirrors meet in a corner at a 90 degree angle.  At what incident angle does the incoming ray leave in a direction parallel to the incoming direction? \\ \vspace{4cm}
\end{enumerate}

\section{Wave Optics}

\begin{enumerate}
\item Suppose a beam of optical light enters glass, and the incident angle with respect to the normal direction at the interface is not equal to the transmitted angle.  If light were a stream of particles, rather than a wave, would that be possible?
\end{enumerate}

\end{document}
